% ============================================================
%   Chapter 6  \textemdash\  Experimentation and Results
% ============================================================
%
% NOTE TO THE AUTHOR: this chapter is written as a STRUCTURAL TEMPLATE.
% Every empirical number appears as "TBD" and is marked with a comment
% explaining the exact source (metrics.sqlite column, dashboard panel,
% etc.). Replace each TBD with the final figures from your runs.
% ============================================================

\chapter{Experimentation and Results}\label{ch:experimentation}

This chapter reports the empirical evaluation of the framework. It is organised around the four
configurations that matter: the baseline agent with no shield (Section~\ref{sec:exp-baseline}),
the basic shield (Section~\ref{sec:exp-basic}), the adaptive-horizon shield
(Section~\ref{sec:exp-adaptive}), and the effect of enabling or disabling the curriculum
(Section~\ref{sec:exp-curriculum}). Section~\ref{sec:exp-protocol} fixes the experimental
protocol and Section~\ref{sec:exp-metrics} formally defines every metric reported hereafter;
Section~\ref{sec:exp-discussion} offers a cross-configuration discussion.

All numerical entries marked as \emph{TBD} are to be filled in with the exact values obtained
from the SQLite metrics database of each run. The SQL queries used to compute the reported
statistics are documented in Appendix~\ref{app:reproducibility}.

\section{Experimental protocol}\label{sec:exp-protocol}

\paragraph{Hardware and software.}
All experiments are carried out on the workstation described in
Section~\ref{sec:req-constraints}:
\begin{itemize}
  \item CPU: \emph{TBD}.% replace with e.g. "AMD Ryzen 7 7800X3D, 8 cores / 16 threads"
  \item GPU: \emph{TBD}.% e.g. "NVIDIA RTX 4070, 12 GB VRAM"
  \item System RAM: \emph{TBD}.% e.g. "32 GB DDR5"
  \item Operating system: Windows 11 Pro.
  \item Python: \emph{TBD}.% e.g. "3.11.9"
  \item PyTorch: \emph{TBD} with CUDA \emph{TBD}.% e.g. "2.3.0 with CUDA 12.1"
  \item CARLA: 0.9.16, launched with \texttt{.\textbackslash CarlaUE4.exe -RenderOffScreen
        -carla-port=2000}.
\end{itemize}

\paragraph{Benchmark scenario.}
The main benchmark is the Town04 highway loop with clear weather (\texttt{ClearNoon}). The
vehicle spawns at a random waypoint each episode and the goal is to travel at least
\(250\,\mathrm{m}\) along the map's route graph within \(1000\) simulator steps
(\(50\,\mathrm{s}\)).

\paragraph{Training budget.}
Every configuration is trained for \(T = 2\,500\) episodes, with PPO updates every
\(2\,048\) steps (\(\sim\!\emph{TBD}\) updates per run).% fill in updates-per-run
The learning rate starts at \(10^{-4}\) and anneals cosine-style to \(10^{-5}\); the target KL
for the early stop is \(0.08\); the entropy bonus is \(0.02\); the value-loss coefficient is
\(0.25\); all other hyperparameters take the defaults of Appendix~\ref{app:hyperparameters}.

\paragraph{Random seeds.}
Every training run uses \(\mathrm{seed}=42\); every evaluation uses \(\mathrm{seed}=100\). Three
independent runs per configuration are carried out when time permits, by varying the seed to
\(\{42, 43, 44\}\); otherwise the single \(\mathrm{seed}=42\) run is used. The reported aggregate
statistics (mean, standard deviation, median) are computed across these three runs when
applicable.

\paragraph{Curriculum.}
Unless explicitly stated otherwise, the curriculum is enabled with \(N_{\max} = 40\). The ablation
of Section~\ref{sec:exp-curriculum} disables it.

\paragraph{Evaluation protocol.}
After training, each agent is evaluated on \(10\) episodes with \(\mathrm{seed}=100\). The shield
configuration used at evaluation is identical to the one used during training: removing the
shield at evaluation would compromise the comparability of the results because a policy trained
with a shield has learned to rely on its interventions.

\section{Metrics}\label{sec:exp-metrics}

The metrics reported throughout this chapter are exactly those produced by the training and
evaluation pipelines. Their operational definitions are collected here to avoid ambiguity.

\begin{description}
  \item[\textbf{Reward/Raw\_Episode.}] Undiscounted sum of shaped rewards of a single episode.
  \item[\textbf{Reward/Average\_100\_Episodes.}] Rolling mean of \emph{Reward/Raw\_Episode} over
        the last \(100\) episodes.
  \item[\textbf{Training/Success\_Rate.}] Fraction of the last \(100\) episodes that ended with
        \texttt{info["arrive\_dest"] = True}.
  \item[\textbf{Training/Crash\_Rate.}] Fraction of the last \(100\) episodes that ended with
        \texttt{info["collision"] = True}.
  \item[\textbf{Training/Offroad\_Rate.}] Fraction of the last \(100\) episodes that ended with
        \texttt{info["out\_of\_road"] = True}.
  \item[\textbf{Safety/Shield\_Activations.}] Number of steps in which the shield overrode the
        agent during the episode.
  \item[\textbf{Safety/Shield\_Rate.}] \emph{Safety/Shield\_Activations} divided by the number of
        steps in the episode.
  \item[\textbf{Safety/Semantic/\{Dynamic,Static,Pedestrian\}\_Interventions.}] Per-episode
        intervention count by threat category, as produced by
        \texttt{CarlaAdaptiveHorizonShield.get\_statistics}.
  \item[\textbf{Safety/Semantic/\{Safe,Warning,Critical\}\_Step\_Rate.}] Fractions of steps spent
        in each risk level.
  \item[\textbf{CARLA/Mean\_Speed\_kmh.}] Arithmetic mean of \texttt{speed\_kmh} across the steps
        of the episode.
  \item[\textbf{CARLA/Speed\_Compliance\_Rate.}] Fraction of steps (with a positive speed limit)
        that satisfy \(\texttt{speed\_kmh} \le 1.05\,\texttt{speed\_limit\_kmh}\).
  \item[\textbf{Safety/Min\_\{Vehicle,Pedestrian\}\_Distance\_m.}] Minimum distance over the
        episode from the ego vehicle to the nearest vehicle / pedestrian.
  \item[\textbf{Loss/Policy\_Loss, Loss/Value\_Loss, Loss/Grad\_Norm, Training/Entropy,
        Training/Approx\_KL, Training/Epochs\_Run.}] Per-update diagnostics produced by the PPO
        loop.
\end{description}

\section{Baseline\,: \texttt{shield\_type=none}}\label{sec:exp-baseline}

The baseline corresponds to an unconstrained PPO agent that follows exactly the same training
protocol but without any shield. Its purpose is twofold: (i) it quantifies what a raw
RL agent can learn on Town04 given the curriculum and the dense reward, and (ii) it provides
the reference numbers against which every shielded configuration is compared.

\paragraph{Training curves.}
Figure~\ref{fig:baseline-reward} shows the rolling reward, the success rate and the safety
rates over the \(2\,500\) training episodes. The agent reaches a final rolling reward of
\emph{TBD} with a success rate of \emph{TBD} and a crash rate of \emph{TBD}.% pull from metrics_episode

\begin{figure}[h]
  \centering
  \fbox{\parbox{0.9\textwidth}{\centering\vspace{1cm} \emph{TBD: figure file, e.g.}
        \texttt{figures/baseline\_reward.pdf} \vspace{1cm}}}
  \caption{Baseline run (\texttt{--shield\_type none}). Top: rolling reward over \(100\)
           episodes. Bottom left: success, crash and off-road rates. Bottom right: number of
           shield activations per episode (identically zero for this configuration).}
  \label{fig:baseline-reward}
\end{figure}

\paragraph{Curriculum trajectory.}
Figure~\ref{fig:baseline-curriculum} shows the evolution of the curriculum stage. Rollbacks, if
any, are visible as downward steps.

\begin{figure}[h]
  \centering
  \fbox{\parbox{0.9\textwidth}{\centering\vspace{1cm} \emph{TBD: figure file, e.g.}
        \texttt{figures/baseline\_curriculum.pdf}\vspace{1cm}}}
  \caption{Baseline run. Curriculum stage (0 to 4) across training episodes and associated NPC
           count.}
  \label{fig:baseline-curriculum}
\end{figure}

\paragraph{Evaluation.}
Table~\ref{tab:baseline-eval} summarises the \(10\) evaluation episodes.

\begin{table}[h]
  \centering
  \caption{Baseline evaluation, \(10\) episodes, \(\mathrm{seed}=100\).}
  \label{tab:baseline-eval}
  \small
  \begin{tabular}{@{}lc@{}}
    \toprule
    \textbf{Metric} & \textbf{Baseline} \\
    \midrule
    Average reward        & TBD $\pm$ TBD \\
    Success rate          & TBD \\
    Crash rate            & TBD \\
    Off-road rate         & TBD \\
    Timeout rate          & TBD \\
    Mean speed (km/h)     & TBD \\
    Speed-compliance rate & TBD \\
    \bottomrule
  \end{tabular}
\end{table}

\section{Basic shield}\label{sec:exp-basic}

The second configuration enables the basic shield (\texttt{--shield\_type basic}). Shield
interventions are now logged per step, per episode and per semantic reason
(\texttt{front}, \texttt{side\_right}, \texttt{side\_left}, \texttt{lane\_right},
\texttt{lane\_left}, \texttt{heading}).

\paragraph{Training curves.}
Figure~\ref{fig:basic-training} mirrors Figure~\ref{fig:baseline-reward}; the additional panel
reports the shield-activation rate per episode. The final rolling reward is \emph{TBD};
the final crash rate is \emph{TBD}, compared to \emph{TBD} for the
baseline.

\begin{figure}[h]
  \centering
  \fbox{\parbox{0.9\textwidth}{\centering\vspace{1cm} \emph{TBD: figure file, e.g.}
        \texttt{figures/basic\_training.pdf}\vspace{1cm}}}
  \caption{Basic-shield run. Top: rolling reward. Bottom left: safety rates.
           Bottom right: shield-activation rate per episode.}
  \label{fig:basic-training}
\end{figure}

\paragraph{Intervention breakdown.}
Table~\ref{tab:basic-breakdown} gives the mean number of interventions per episode by sub-reason,
averaged over the last \(100\) training episodes.

\begin{table}[h]
  \centering
  \caption{Basic shield: interventions per episode by sub-reason (last \(100\) episodes).}
  \label{tab:basic-breakdown}
  \small
  \begin{tabular}{@{}lc@{}}
    \toprule
    \textbf{Sub-reason}  & \textbf{Mean interventions / episode} \\
    \midrule
    \texttt{front}       & TBD \\
    \texttt{side\_right} & TBD \\
    \texttt{side\_left}  & TBD \\
    \texttt{lane\_right} & TBD \\
    \texttt{lane\_left}  & TBD \\
    \texttt{heading}     & TBD \\
    \bottomrule
  \end{tabular}
\end{table}

\paragraph{Evaluation.}
Table~\ref{tab:basic-eval} reports the \(10\) evaluation episodes.

\begin{table}[h]
  \centering
  \caption{Basic shield evaluation, \(10\) episodes, \(\mathrm{seed}=100\).}
  \label{tab:basic-eval}
  \small
  \begin{tabular}{@{}lc@{}}
    \toprule
    \textbf{Metric} & \textbf{Basic shield} \\
    \midrule
    Average reward           & TBD $\pm$ TBD \\
    Success rate             & TBD \\
    Crash rate               & TBD \\
    Off-road rate            & TBD \\
    Timeout rate             & TBD \\
    Mean speed (km/h)        & TBD \\
    Speed-compliance rate    & TBD \\
    Shield activations / ep  & TBD \\
    \bottomrule
  \end{tabular}
\end{table}

\section{Adaptive-horizon shield}\label{sec:exp-adaptive}

The third configuration enables the adaptive-horizon shield
(\texttt{--shield\_type adaptive}). In addition to the metrics reported for the basic shield,
the adaptive shield exposes the time spent in each risk level and the breakdown of
interventions by threat category (dynamic, static, pedestrian).

\paragraph{Training curves.}
Figure~\ref{fig:adaptive-training} mirrors Figure~\ref{fig:basic-training} with an extra panel
for the risk-level time fractions.

\begin{figure}[h]
  \centering
  \fbox{\parbox{0.9\textwidth}{\centering\vspace{1cm} \emph{TBD: figure file, e.g.}
        \texttt{figures/adaptive\_training.pdf}\vspace{1cm}}}
  \caption{Adaptive-shield run. Top-left: rolling reward. Top-right: safety rates.
           Bottom-left: shield-activation rate by threat category. Bottom-right: fraction of
           steps per risk level.}
  \label{fig:adaptive-training}
\end{figure}

\paragraph{Risk distribution.}
Table~\ref{tab:adaptive-risk} reports the average time spent in each risk level over the last
\(100\) training episodes.

\begin{table}[h]
  \centering
  \caption{Adaptive shield: average time per risk level (last \(100\) episodes).}
  \label{tab:adaptive-risk}
  \small
  \begin{tabular}{@{}lc@{}}
    \toprule
    \textbf{Risk level}  & \textbf{Fraction of steps} \\
    \midrule
    safe     & TBD \\
    warning  & TBD \\
    critical & TBD \\
    \bottomrule
  \end{tabular}
\end{table}

\paragraph{Intervention breakdown.}
Table~\ref{tab:adaptive-breakdown} reports the intervention count per episode by threat
category.

\begin{table}[h]
  \centering
  \caption{Adaptive shield: interventions per episode by threat category (last \(100\)
           episodes).}
  \label{tab:adaptive-breakdown}
  \small
  \begin{tabular}{@{}lc@{}}
    \toprule
    \textbf{Category}    & \textbf{Mean interventions / episode} \\
    \midrule
    dynamic    & TBD \\
    static     & TBD \\
    pedestrian & TBD \\
    \bottomrule
  \end{tabular}
\end{table}

\paragraph{Evaluation.}
Table~\ref{tab:adaptive-eval} reports the \(10\) evaluation episodes.

\begin{table}[h]
  \centering
  \caption{Adaptive shield evaluation, \(10\) episodes, \(\mathrm{seed}=100\).}
  \label{tab:adaptive-eval}
  \small
  \begin{tabular}{@{}lc@{}}
    \toprule
    \textbf{Metric} & \textbf{Adaptive shield} \\
    \midrule
    Average reward                   & TBD $\pm$ TBD \\
    Success rate                     & TBD \\
    Crash rate                       & TBD \\
    Off-road rate                    & TBD \\
    Timeout rate                     & TBD \\
    Mean speed (km/h)                & TBD \\
    Speed-compliance rate            & TBD \\
    Shield activations / ep          & TBD \\
    Mean time in \emph{critical} (\%)& TBD \\
    \bottomrule
  \end{tabular}
\end{table}

\paragraph{Cross-configuration comparison.}
Table~\ref{tab:config-comparison} aggregates the evaluation of the three configurations
side-by-side.

\begin{table}[h]
  \centering
  \caption{Side-by-side evaluation comparison (\(10\) episodes, seed \(100\)).}
  \label{tab:config-comparison}
  \small
  \begin{tabular}{@{}lccc@{}}
    \toprule
    \textbf{Metric}            & \textbf{None} & \textbf{Basic} & \textbf{Adaptive} \\
    \midrule
    Average reward             & TBD & TBD & TBD \\
    Success rate               & TBD & TBD & TBD \\
    Crash rate                 & TBD & TBD & TBD \\
    Off-road rate              & TBD & TBD & TBD \\
    Speed-compliance rate      & TBD & TBD & TBD \\
    Min. vehicle distance (m)  & TBD & TBD & TBD \\
    Shield activations / ep    & 0   & TBD & TBD \\
    \bottomrule
  \end{tabular}
\end{table}

\section{Effect of the curriculum}\label{sec:exp-curriculum}

To quantify the contribution of the curriculum we repeat the adaptive-shield configuration with
\texttt{--no-curriculum}, i.e.\ with the manager disabled and the agent exposed to the maximum
NPC count from episode~\(1\). Everything else is held fixed.

\paragraph{Training curves.}
Figure~\ref{fig:curriculum-ablation} compares the rolling reward of both runs.
Without the curriculum, the agent takes \emph{TBD} episodes to reach the reward level that the
curriculum-enabled run reaches after \emph{TBD} episodes; the final rolling crash rate without
curriculum is \emph{TBD} against \emph{TBD} with curriculum.% pull from metrics_episode

\begin{figure}[h]
  \centering
  \fbox{\parbox{0.9\textwidth}{\centering\vspace{1cm} \emph{TBD: figure file, e.g.}
        \texttt{figures/curriculum\_ablation.pdf}\vspace{1cm}}}
  \caption{Curriculum ablation: adaptive shield with curriculum enabled (solid) vs.\ disabled
           (dashed). Top: rolling reward. Bottom: rolling crash rate.}
  \label{fig:curriculum-ablation}
\end{figure}

\paragraph{Rollback analysis.}
When the curriculum is enabled, Table~\ref{tab:curriculum-events} enumerates the stage
transitions and rollbacks recorded during training.

\begin{table}[h]
  \centering
  \caption{Curriculum events during the adaptive-shield run.}
  \label{tab:curriculum-events}
  \small
  \begin{tabular}{@{}llll@{}}
    \toprule
    \textbf{Episode} & \textbf{Event} & \textbf{From stage} & \textbf{To stage} \\
    \midrule
    TBD & advance  & TBD & TBD \\
    TBD & advance  & TBD & TBD \\
    TBD & rollback & TBD & TBD \\
    \emph{\ldots} & \emph{\ldots} & \emph{\ldots} & \emph{\ldots} \\
    \bottomrule
  \end{tabular}
\end{table}

\section{Discussion}\label{sec:exp-discussion}

This section will discuss the trade-offs observed across configurations. The analysis will cover
at least the following four axes; each is introduced here as a question that the empirical
numbers of Sections~\ref{sec:exp-baseline}--\ref{sec:exp-curriculum} will answer.

\begin{enumerate}
  \item \textbf{Safety vs.\ performance.} How much reward does a shield cost in exchange for a
        lower crash rate? Is the adaptive shield strictly Pareto-better than the basic shield,
        or does one of them dominate on one axis while the other dominates on the complementary
        axis?
  \item \textbf{Intervention regime.} What fraction of a typical episode is spent with the
        shield active, and how does this fraction evolve during training? A decreasing
        shield-activation rate is consistent with the hypothesis that the policy is internalising
        the shield's behaviour.
  \item \textbf{Threat composition.} Which fraction of interventions is due to dynamic threats
        (other vehicles), static threats (barriers, walls) or pedestrians? How do these
        fractions change across curriculum stages?
  \item \textbf{Curriculum value.} Does the performance-driven curriculum accelerate learning
        and, more importantly, does it reduce the worst-case transient crash rate of the
        adaptive-shield run? The empirical answer to this question informs the recommendation
        of whether to enable the curriculum by default.
\end{enumerate}

Each enumerated item will be supported by a specific cross-reference to the tables and figures
above. Once the TBD entries have been filled in, a concise summary (bulleted or tabular) will be
inserted at the end of this section to provide a reader-friendly overview of the
findings.
